\documentclass[../Main.tex]{subfiles}

\begin{document}

\section{Verzeichnisse}

\subsection{Literatur- und Quellenverzeichnis}

\printbibliography[heading=none, nottype=software]

\subsection{Ludografie}

\printbibliography[heading=none, type=software]

\subsection{Abbildungsverzeichnis}

\listoffigures

\subsection{Tabellenverzeichnis}

\listoftables

\subsection{Beteiligte Personen} \label{sec:Beteiligte}

An der Entstehung dieser Arbeit waren die folgenden Personen beteiligt.\footnote{Mit Claude formuliert, siehe KI-Deklaration (Abschnitt~\ref{sec:KIDeklaration}).} Ihnen allen danken wir herzlich für ihre Zeit und ihre Unterstützung.

\begin{description}[leftmargin=1.5cm, style=nextline, itemsep=0.6em]
    \item[Projektteam] Rodin Arslan (\url{rodin.delil.arslan@gmail.com}) und Zhile Zhu (\url{zhuzhile08@gmail.com}), Klasse G23g, Kantonsschule Baden: Konzeption, Entwicklung des Prototyps und Verfassen des Begleitberichts
    \item[Erstbetreuung] Michael Schneider, Kantonsschule Baden, Dept.\ Informatik, der uns während der gesamten Arbeit begleitet hat
    \item[Zweitbewerter] Dr.\ Christophe Bersier, Kantonsschule Baden, Dept.\ Physik
    \item[Testpersonen] Bastian Nakatsu, Beatrice Schötzau, Cyrill Bruttin, Luca Savoldelli und Max Starrowsky, die den Prototyp in den Playtests gespielt und uns mit ihren Rückmeldungen wertvolle Hinweise für unsere Auswertung gegeben haben (Abschnitt~\ref{sec:Durchfuehrung}). In den Protokollen (Abschnitt~\ref{sec:PlaytestProtokolle}) sind sie anonymisiert. Die alphabetische Reihenfolge hier entspricht nicht der Nummerierung TP1--TP5.
    \item[Weitere Unterstützung] Beatrice Schötzau (\url{beatrice.schotzau@gmail.com}) hat das coole App-Icon des Prototyps gestaltet.
\end{description}

\subsection{KI-Deklaration} \label{sec:KIDeklaration}

Für diese Arbeit haben wir KI-Werkzeuge eingesetzt: beim Begleitbericht und bei einzelnen Teilen des Prototyps Claude (Anthropic). Die folgende Tabelle fasst den Einsatz nach den Bereichen des Rasters \enquote{Verwendung von KI} zusammen. Die einzelnen Nachweise mit Prompt, Datum und Art der Übernahme finden sich in unseren Arbeitsjournalen. An Stellen, an denen von der KI formulierte Texte in den Begleitbericht eingeflossen sind, steht zusätzlich eine Fussnote. Diese Nachweise sind am Ende dieses Abschnitts aufgeführt. \textcolor{red}{Bei Programmcode ist der KI-Einsatz einschliesslich der verwendeten Prompts direkt im Quellcode der betroffenen Dateien dokumentiert.}

KI haben wir nicht als Quelle für sachliche Informationen verwendet. Alle Sachaussagen stützen sich auf die angegebenen Quellen. Quellenangaben, welche die KI geprüft oder formatiert hat, haben wir selbst mit den Originalquellen abgeglichen.

{\small
    \renewcommand{\arraystretch}{1.3}
    \begin{longtable}{|>{\raggedright\arraybackslash}p{2.3cm}|>{\raggedright\arraybackslash}p{2.4cm}|>{\raggedright\arraybackslash}p{4.4cm}|>{\raggedright\arraybackslash}p{4.6cm}|}
        \hline
        \textbf{Bereich} & \textbf{Tool, Datum} & \textbf{Zweck und betroffene Stelle} & \textbf{Art der Übernahme, Eigenanteil und Überprüfung} \\
        \hline
        \endfirsthead
        \hline
        \textbf{Bereich} & \textbf{Tool, Datum} & \textbf{Zweck und betroffene Stelle} & \textbf{Art der Übernahme, Eigenanteil und Überprüfung} \\
        \hline
        \endhead
        \multicolumn{4}{r}{\textit{Fortsetzung auf der nächsten Seite}} \\
        \endfoot
        \caption{Einsatz von KI nach den Bereichen des Rasters \enquote{Verwendung von KI}}
        \label{tab:KI}
        \endlastfoot
        Arbeitsplanung und Organisation & Claude Opus 5.5 (Claude Cowork), 28.10.--9.11.2026 & Aufgabenliste für den Begleitbericht, Abschätzen des verbleibenden Aufwands & Liste als Orientierung verwendet und laufend selbst angepasst. \\
        \hline
        Informations- und Literaturrecherche & Claude (claude.ai), Juli 2026; Claude Opus 5.5, 28.10.--4.11.2026 & Vorschläge für Fachbücher zum Leveldesign auf unsere Frage nach Material; BibTeX-Einträge zu unseren Quellen, Einschätzung ihrer Eignung, Prüfen von Quellenangaben (Literaturverzeichnis) & Vorgeschlagene Bücher selbst gelesen und geprüft. Die übrigen Quellen haben wir selbst über Google Scholar gesucht und gelesen. BibTeX-Einträge übernommen und mit den Originalquellen abgeglichen. \\
        \hline
        Vorgehen und Methode & Claude Opus 5.5, 28.10.--9.11.2026 & Rückmeldungen zur Beschreibung des Vorgehens (Kapitel~\ref{sec:Grundlagen}), Anpassen der Beschreibung des Beobachtungsrasters an das tatsächliche Vorgehen & Hinweise geprüft und Texte überwiegend selbst überarbeitet. MDA-Fragebogen und Playtest-Fragebögen haben wir selbst entwickelt. \\
        \hline
        Untersuchung und Datenerhebung & -- & Kein KI-Einsatz: Die Playtests haben wir selbst geplant und durchgeführt. & -- \\
        \hline
        Aufbereitung von Daten & Claude Opus 5.5, 28.10.--9.11.2026 & Zuordnen der Beobachtungsnotizen zu den Kriterien B1--B14, sprachliches Vereinheitlichen der Notizen und Interviewantworten, Anonymisieren der Protokolle (Anhang~\ref{sec:PlaytestProtokolle}); Heraussuchen der Stellen aus den rund 20 Seiten Protokollen, die für die einzelnen Designvorschläge passen, und Auszählen der Antworten für Kapitel~\ref{sec:Evaluation} (z.\,B.\ \enquote{x von 5}) & Zuordnung und Formulierungen weitgehend übernommen, von uns geprüft und an mehreren Stellen korrigiert. Die herausgesuchten Stellen haben wir anhand der Protokolle geprüft und ergänzt, damit nichts fehlt. Zählungen ebenfalls anhand der Protokolle geprüft. Zu den Interpretationen in Kapitel~\ref{sec:Evaluation} siehe die Zeile \enquote{Weiteres: von der KI formulierte Texte}. \\
        \hline
        Hilfsmittel und Werkzeuge & Claude Opus 5.5, 28.10.--9.11.2026 & Textsatz mit \LaTeX: mehrseitige Tabellen, Tabellenformatierung, Fehlersuche beim Kompilieren & \LaTeX-Code übernommen und im kompilierten Dokument kontrolliert. \\
        \hline
        Programmieren und Berechnungen & Claude Sonnet 4.6 (claude.ai), 04.03.2026 & Editor-Werkzeug zum Erzeugen von Treppen & Code übernommen; Prompts \textcolor{red}{im Quellcode dokumentiert}. \\
        \hline
        & Claude Sonnet 4.6 (claude.ai), 09.04.2026 und 22.05.2026 & Hilfsskript für Animationen, Grundgerüst und eine Funktion & Als Ausgangspunkt übernommen und danach von uns verändert; Prompts \textcolor{red}{im Quellcode dokumentiert}. \\
        \hline
        & Claude Sonnet 5.5 (claude.ai), Oktober 2026 & Audiosystem: Musikverwaltung und Grundlage für Soundeffekte & Grösstenteils übernommen; Prompts \textcolor{red}{im Quellcode dokumentiert}. \\
        \hline
        & Claude bzw.\ ChatGPT (jeweils Gratisversion), Juli und Oktober 2026 & \textcolor{red}{Shader: nicht kompatible GLSL-Shader nach Godot portiert (Juli) sowie zwei eigene Stilisierungs-Shader für die Toon-Beleuchtung und die Nachbearbeitung des Bildes geschrieben (Oktober)} & Grösstenteils übernommen und im Spiel visuell überprüft. \\
        \hline
        & Claude bzw.\ ChatGPT (jeweils Gratisversion), 17.10.2026 & \textcolor{red}{Skript, das die Hitboxen der Blender-Modelle automatisch erzeugt} & \textcolor{red}{Übernommen und an den Modellen im Spiel überprüft.} \\
        \hline
        Rückmeldungen und Orientierung & Claude Opus 5.5, 28.10.--9.11.2026 & Durchsicht aller Kapitel des Begleitberichts mit Hinweisen zu Struktur, Verständlichkeit, fehlenden Belegen und Widersprüchen zwischen den Kapiteln & Hinweise als Notizen erhalten. Die Texte haben wir überwiegend selbst überarbeitet, einige Hinweise bewusst verworfen. \\
        \hline
        Überprüfung und Kontrolle & Claude Opus 5.5, 28.10.--9.11.2026 & Korrektur von Rechtschreibung, Grammatik und Wortwahl; Kennzeichnen von eigenen und belegten Aussagen; Abgleich mit Projektvertrag und Zwischenbewertung; Prüfen der Einheitlichkeit von Begriffen und Nummerierungen & Korrekturen übernommen. Der Inhalt wurde dabei nicht verändert. \\
        \hline
        Darstellung von Ergebnissen & Claude Opus 5.5, 28.10.--9.11.2026 & Gestaltung der Tabellen~\ref{tab:Designvorschlaege} bis~\ref{tab:Raeume} aus unseren Angaben & Aufbau der Tabellen übernommen. Die Inhalte stammen von uns. \\
        \hline
        Weiteres: von der KI formulierte Texte & Claude Opus 5.5, 28.10.--9.11.2026 & Glossar (Anhang~\ref{sec:Glossar}), Abschnitt~\ref{sec:Beteiligte} (Beteiligte Personen) und Abschnitt~\ref{sec:KIDeklaration} (KI-Deklaration) & Aus unseren Angaben formuliert und weitgehend übernommen. Inhaltlich von uns geprüft. An den Stellen mit Fussnote gekennzeichnet. \\
        \hline
        Weiteres: Arbeitsjournal & Claude bzw.\ ChatGPT (jeweils Gratisversion), 29.03.--10.11.2026 & Strukturieren unserer Journaleinträge nach den Vorgaben der Betreuungsperson & Gliederung übernommen. Die Inhalte stammen von uns. \\
        \hline
    \end{longtable}
}

Nicht mit KI erstellt haben wir insbesondere die Analysen der Referenzspiele, die Designvorschläge, den Kampfabschnitt des Prototyps mit seinen Räumen und Gegnern\textcolor{red}{ und die Durchführung der Playtests. Bei der Interpretation der Ergebnisse hat die KI fehlende Belege, Urteile und Vergleiche vorgeschlagen. Entschieden haben wir darüber selbst.}

\textcolor{red}{Die KI-Einsätze von Zhile Zhu, etwa beim Programmcode und beim Projektvertrag, sind in seinem eigenen Arbeitsjournal nachgewiesen.}

\paragraph{Nachweise der von der KI formulierten Texte}
\begin{itemize}[leftmargin=1.5cm, itemsep=0.4em]
    \item Claude (Opus 5.5): «[\ldots] leg ein glossar am ende an [\ldots] und auch bitte eine asset liste», 01.11.2026. Glossar nach unseren Vorgaben verfasst und weitgehend übernommen.
    \item Claude (Opus 5.5): «die limitationen von unserem playtest bogen kannst du aufschreiben [\ldots]», 01.11.2026. Abschnitt Limitationen aus unseren Stichpunkten verfasst und weitgehend übernommen.
    \item Claude (Opus 5.5): «also ich habe jetzt v6 geschreiebn. das mit den hindernissen und den verschiedenen ebenen kannst du schreiben und auch das mit auf die karte verweisen. schau dir hier besonders die struktur an ich glaube es ist ein wenig alles überll und wenig struktur», 04.11.2026. Absatz zu Hindernissen und Ebenen in Abschnitt~\ref{sec:V6} formuliert und übernommen. Ebenso einzelne Rückbezüge in den Abschnitten~\ref{sec:Kampfabschnitt} bis~\ref{sec:V8}.
    \item Claude (Opus 5.5): «okay ich habe jetzt auch den letzten palytest gemacht. damit bin ich mit allen fertig. bitte korrigiere noch testperson eins und füge alle sachen an ihren platzt [\ldots]», 04.11.2026. Beobachtungsnotizen den Kriterien zugeordnet und sprachlich vereinheitlicht; übernommen und korrigiert.
    \item Claude (Opus 5.5): «okay jetzt was ich jetzt von dir will ist dass du die die tests komplett anonymisierst und auch das geschriebene genderneutral aufschreibst, und [\ldots] alle personen die beteiligt sind an unserem projekt auf eine seite schreiben [\ldots]», 04.11.2026. Protokolle anonymisiert und Abschnitt Beteiligte Personen verfasst; übernommen.
    \item \textcolor{red}{Claude (Opus 5.5): «okay setze alle punkte in kapitel 5 um und lasse was du eingefügt hast rot dass ichs korrigieren/ändern kann und in kapitel 6, [\ldots] setze auch hier alle deine vorschläge um», 09.11.2026, sowie gleichlautende Anweisungen zu V3--V8 und zum Gesamteindruck am selben Tag. Vorschläge zu fehlenden Belegen, Urteilen, geplanten und ungeplanten Wirkungen und Vergleichen mit den Referenzspielen in den Abschnitten~\ref{sec:V1} bis~\ref{sec:V8} und~\ref{sec:ErgVorschlag} formuliert. Von uns einzeln geprüft, teils geändert und übernommen.}
    \item \textcolor{red}{Claude (Opus 5.5): «[\ldots] Du musst mir einfach jetzt diese, wie viele waren es? 5 Artikel schreiben. [\ldots] Schreib mir einfach die Reflexionen [\ldots]», 08.11.2026. Arbeitsjournal-Einträge vom 26.10. bis 08.11.2026 aus unseren Angaben formuliert. Von uns geprüft und ergänzt.}
\end{itemize}

\begin{itemize}[leftmargin=1.5cm]
    \item \textcolor{red}{\textbf{Zitierregeln:} KI-Ergebnisse dürfen nicht wortwörtlich übernommen werden (Merkblatt, Zitierregeln 1.7). Glossar, Limitationen und Beteiligte Personen entweder in eigenen Worten umschreiben oder die Übernahme mit Herrn Schneider absprechen. Das betrifft jetzt auch die von der KI formulierten Sätze in Kapitel~\ref{sec:Umsetzung} und~\ref{sec:Evaluation}. Danach die Zeile \enquote{Weiteres: von der KI formulierte Texte} anpassen.}
\end{itemize}

\newpage
\addtocontents{toc}{\protect\setcounter{tocdepth}{2}}
\section{Anhang: Playtest-Protokolle} \label{sec:PlaytestProtokolle}

In diesem Anhang sind die Daten der fünf Playtests dokumentiert. Pro Testperson sind die Antworten des Pre-Questionnaire (P1--P8) und des Post-Questionnaire (Q1--Q14), die Notizen zum Beobachtungsraster (B1--B14) sowie die Antworten aus dem semistrukturierten Interview (I1--I8) aufgeführt. Die Nummerierung entspricht derjenigen in Abschnitt~\ref{sec:VorgehenPlaytest}. Die zugehörigen Fragen und Kriterien finden sich in den Abschnitten~\ref{sec:PreQuestionnaire} (Pre-Questionnaire), \ref{sec:PostQuestionnaire} (Post-Questionnaire), \ref{sec:Beobachtungsraster} (Beobachtungsraster) und~\ref{sec:Interview} (Interview). Die Aussagen Q1--Q11 wurden auf einer Skala von 1 (\enquote{trifft gar nicht zu}) bis 5 (\enquote{trifft völlig zu}) bewertet. Da das Interview an die beobachteten Auffälligkeiten angepasst wurde, sind dort jeweils die Fragen aufgeführt, wie sie der Testperson tatsächlich gestellt wurden. Während der Playtests verwendeten wir noch eine ältere Raumnummerierung. Die Raumnummern in den Notizen haben wir an die Nummern der Encounter in Tabelle~\ref{tab:Raeume} angepasst.

Die Testpersonen sind anonymisiert und werden nur mit TP1 bis TP5 bezeichnet. Namen, die in den Antworten vorkamen, sind durch [Name] ersetzt. In den Beobachtungsnotizen sprechen wir unabhängig vom Geschlecht immer von \enquote{der Testperson} bzw.\ von \enquote{sie}. Die Antworten des Pre- und des Post-Questionnaire sind wörtlich übernommen. Die Antworten aus dem Interview geben wir sinngemäss wieder: Rechtschreibung und Zeichensetzung haben wir korrigiert und die Sätze teilweise sprachlich geglättet, umgangssprachliche Ausdrücke aber beibehalten. Die Beobachtungsnotizen haben wir nachträglich den Kriterien B1--B14 zugeordnet und sprachlich vereinheitlicht.\footnote{Zuordnung, sprachliche Vereinheitlichung und Anonymisierung mit Claude, siehe KI-Deklaration (Abschnitt~\ref{sec:KIDeklaration}).}

\subsection{Testperson 1 (TP1)}

\subsubsection{Pre-Questionnaire}

\begin{enumerate}[label=\textbf{P\arabic*}, leftmargin=1.5cm, itemsep=0.5em]
    \item Während der Schulzeit nicht aber in den Ferien etwa die hälfte der Ferien.
    \item Survival und Shooter
    \item Die Grafik, wenn man lange braucht bis man irgendwas versteht und dass man wegrennen kann
    \item Ich stelle mir vor dass man Kämpft und versucht herauszufinden, wie etwas funktioniert.
    \item 3-4
    \item Ganz klar aggressiv und etwas methodisch wenn man das Spiel besser kann
    \item Ds Game kennenlernen, woran ihr jetzt anscheinend so lange gearbeitet habt. Ich habe ja ein paar Ausschnitte bei euch gesehen
    \item Ich finde Unübersichtlichkeit ehrlich gesagt ganz wizig, also es nervt eigentlich richtig, aber es fühlt sich dann belohnend an. Als das schätze ich an einem Game. Zu leichte gegner sind langweilig.
\end{enumerate}

\subsubsection{Post-Questionnaire}

\paragraph{Gegnerverhalten}
\begin{enumerate}[label=\textbf{Q\arabic*}, leftmargin=1.5cm, itemsep=0.5em]
    \item 3
    \item 2
    \item 4
    \item 2
\end{enumerate}

\paragraph{Gegnerplatzierung und Kampfraum}
\begin{enumerate}[label=\textbf{Q\arabic*}, leftmargin=1.5cm, itemsep=0.5em, resume]
    \item 4
    \item 4
    \item 5
    \item 5
\end{enumerate}

\paragraph{Gesamteindruck}
\begin{enumerate}[label=\textbf{Q\arabic*}, leftmargin=1.5cm, itemsep=0.5em, resume]
    \item 4
    \item 5
    \item 5
    \item der moment wo ich so low healt weglaufen konnte, und der schluss wo ich einmal gestorben bin, die combos, die ich hatte waren cool
    \item ich wollte immer aus der luft auf den gegner schlagen und dann war ich in der luft stuck, habe leider keine bugs gefunden. und ich hatte gehofft, dass es einen endboss gibt.
    \item so habe etwa das erwartet.
\end{enumerate}

\subsubsection{Beobachtungsraster}

\paragraph{Gegnerverhalten}
\begin{enumerate}[label=\textbf{B\arabic*}, leftmargin=1.5cm, itemsep=0.8em]
    \item
        \begin{itemize}[nosep]
            \item Wenn die Testperson zu viel Schaden nimmt, springt sie aus dem Kampf heraus.
        \end{itemize}
    \item
        \begin{itemize}[nosep]
            \item Um eine Welle etwas zu zerstreuen, rennt sie hin und her.
            \item Wenn sie nicht einschätzen kann, wie viele Gegner da sind, rennt sie weg.
        \end{itemize}
    \item
        \begin{itemize}[nosep]
            \item Sie versucht, zu den Archers zu gelangen, wird aber von den Gegnern dazwischen abgehalten.
            \item Sie geht zuerst auf die Archers, die am höchsten platziert sind.
        \end{itemize}
    \item
        \begin{itemize}[nosep]
            \item Beim Springen auf die Boxen wird sie von den Archers in der Luft abgeschossen.
            \item Sie versucht zwar auszuweichen, ignoriert die Angriffe aber oft einfach.
        \end{itemize}
    \item
        \begin{itemize}[nosep]
            \item Sie kann die Angriffsdistanz nicht ganz einschätzen. Auch später schätzt sie die Distanzen noch nicht richtig ein.
            \item Sie hat Mühe, von der Bewegung auf das Kämpfen umzusteigen, und verfehlt dadurch andauernd die Archers.
            \item Sie spammt die Sprungtaste, statt sie gedrückt zu halten.
        \end{itemize}
\end{enumerate}

\paragraph{Gegnerplatzierung und Kampfraum}
\begin{enumerate}[label=\textbf{B\arabic*}, leftmargin=1.5cm, itemsep=0.8em, resume]
    \item
        \begin{itemize}[nosep]
            \item Sie greift einfach an, was gerade vor ihr liegt.
            \item Wenn eine neue Welle erscheint, springt sie heraus und nimmt Abstand.
            \item Sie sagt die Platzierung der Gegner auch in neuen Räumen voraus, in denen sie noch nie war.
            \item Heavy Soldiers drängen sie in eine Ecke, aus der sie nicht mehr herauskommt.
        \end{itemize}
    \item
        \begin{itemize}[nosep]
            \item Keine Einträge notiert.
        \end{itemize}
    \item
        \begin{itemize}[nosep]
            \item Um Abstand vom Kampf zu nehmen, geht sie auf die Boxen.
            \item Sie nimmt hinter Boxen Deckung, um sich zu heilen und vor den Schüssen der Archers geschützt zu sein.
            \item Sie achtet darauf, dass die Energieleiste nie voll ist, und heilt sich, sobald sie voll ist, auch wenn ihr nur wenig Leben fehlt.
            \item Sie heilt sich regelmässig, wenn sie kann.
            \item Manchmal vergisst sie, sich vor dem Raumwechsel zu heilen, und heilt sich dann direkt nach dem Betreten des nächsten Raumes.
        \end{itemize}
    \item
        \begin{itemize}[nosep]
            \item Als sie die Archers zum ersten Mal sieht, greift sie diese nicht an, sondern weicht ihren Schüssen aus.
            \item Beim ersten Heavy Soldier wirkt sie schockiert.
            \item Sie bemerkt sofort, dass der Heavy Soldier nicht genau trifft.
            \item Sie bemerkt sofort, dass der Heavy Archer weniger Leben hat als der normale Archer und dass seine Schüsse anders sind.
            \item Den neuen Chain Archer erkennt sie an dessen Angriffen.
            \item Neue Gegner erkennt sie schon von Weitem an der Farbe.
        \end{itemize}
\end{enumerate}

\paragraph{Allgemein}
\begin{enumerate}[label=\textbf{B\arabic*}, leftmargin=1.5cm, itemsep=0.8em, resume]
    \item
        \begin{itemize}[nosep]
            \item Sie fällt bzw.\ springt die Klippe hinunter und weiss nicht, dass man dabei stirbt.
        \end{itemize}
    \item
        \begin{itemize}[nosep]
            \item Sie fragt, ob man Fallschaden nehmen kann.
            \item In Raum~11 sagt sie voraus, dass sie sterben wird. Das liegt vielleicht an der Spannung, die dieser Raum erzeugt (Vermutung der Beobachtenden).
        \end{itemize}
    \item
        \begin{itemize}[nosep]
            \item Wenn sie die Tasten spammt und die Spielfigur trotzdem ihre Combo ausführt, findet sie das witzig.
        \end{itemize}
    \item
        \begin{itemize}[nosep]
            \item Am Ende einer Welle drückt sie sich gegen die Tür. Es sieht so aus, als würde sie hindurchglitchen.
            \item Später versucht sie erneut, durch die Tür zu glitchen, obwohl die Welle noch läuft.
            \item Sie springt auf Gegner und will zuschlagen, beginnt dabei aber eine Luftcombo und trifft nichts.
            \item Sie versucht, die Gegner dazu zu bringen, sich gegenseitig zu treffen.
            \item Sie stellt sich oft mit dem Rücken an die Wand, damit die Gegner nur von einer Seite kommen, bleibt dann aber stecken.
            \item Sie sucht nach einem Konami-Code.
            \item Sie versucht weiterhin, Glitches zu finden.
            \item Am Ende des Spiels sucht sie nach einem geheimen Raum.
        \end{itemize}
    \item
        \begin{itemize}[nosep]
            \item Sie sieht sofort, wenn sich die Türen öffnen.
            \item Ob ein Encounter fertig ist, kann sie manchmal gut einschätzen, manchmal aber nicht.
            \item Sie bemerkt, dass der Kampfabschnitt ein Rundgang ist.
        \end{itemize}
\end{enumerate}

\subsubsection{Semistrukturiertes Interview}

\begin{enumerate}[label=\textbf{I\arabic*}, leftmargin=1.5cm, itemsep=1em]
    \item
    \begin{enumerate}[label=\alph*), itemsep=0.6em]
        \item Was ist dir durch den Kopf gegangen, als du zum ersten Mal den Heavy Soldier gesehen hast? \\
            \textit{Antwort:} Ich dachte, er nimmt mich richtig hops und dass er schneller schlägt. Er war einfacher als erwartet. \\
            \textit{Bezug zur Beobachtung (B9):} Beim ersten Heavy Soldier wirkte sie schockiert. Sie bemerkte aber sofort, dass der Heavy Soldier nicht genau trifft.
        \item Was ist dir durch den Kopf gegangen, als du die ganze Zeit in der Luft von den Chain Archers abgeschossen wurdest? \\
            \textit{Antwort:} Die haben mich richtig abgefuckt und waren der Grund, wieso ich, glaube ich, auch gestorben bin, weil ich gestunnt war. Ich glaube, die waren insgesamt die schwersten Gegner. \\
            \textit{Bezug zur Beobachtung (B4):} Beim Springen auf die Boxen wurde sie von den Archers in der Luft abgeschossen.
    \end{enumerate}

\item Was hast du beim letzten Raum gedacht, als du hereingekommen bist? \\
    \textit{Antwort:} Ich hatte richtig viele Gegner erwartet, und da war am Anfang nur ein Gegner, aber dann nahm es richtig Fahrt auf.

\item An dieser Stelle bist du wiederholt gestorben bzw.\ hast du gezögert. Was hat dir das Weiterkommen erschwert? \\
    \textit{Antwort:} --

\item Wie hat sich der enge Raum am Schluss anders angefühlt als z.\,B.\ der Brückenraum? \\
    \textit{Antwort:} Im Brückenraum gab es mehr Möglichkeiten, sich zu verstecken und den Gegnern auszuweichen. Stressiger waren sie aber nicht. Die Gegner waren einfach verteilter.

\item Als hier mehrere Gegner gleichzeitig da waren: Wie hast du entschieden, wen du zuerst angreifst? \\
    \textit{Antwort:} Einfach mal probiert, also zuerst den Chain Archer, weil ich den nicht mochte und er stark gegen mich war. Danach die Soldiers auf dem Weg und am Schluss die Heavy Archers. Intuition. \\
    \textit{Bezug zur Beobachtung (B3):} Sie versuchte, zu den Archers zu gelangen, wurde aber von den Gegnern dazwischen abgehalten. Sie griff zuerst die am höchsten platzierten Archers an.

\item Empfandest du diesen Treffer bzw.\ diesen Tod als eigenen Fehler oder als unfair? Warum? \\
    \textit{Antwort:} Ja, am Anfang, bis ich gecheckt habe, was die neuen Gegner machen, aber danach war es gut. \\
    \textit{Bezug zur Beobachtung (B9):} Neue Gegnervarianten erkannte sie schnell an Farbe, Angriffen und Eigenschaften.

\item Gab es einen Punkt, an dem sich der Kampf für dich in einen anderen \enquote{Modus} verschoben hat (z.\,B. von entspannt zu angespannt)? Wo genau? \\
    \textit{Antwort:} Gegen Ende wurde es immer cooler, weil verschiedene Gegner schon eingeführt waren und es verschiedene Kombinationen gab. Da hat es mehr Spass gemacht.

\item Gibt es etwas, das du während des Spielens gedacht oder gefühlt hast, das wir bisher nicht angesprochen haben? \\
    \textit{Antwort:} Wieso war ich eigentlich immer in einem geschlossenen Raum?
\end{enumerate}

\subsection{Testperson 2 (TP2)}

\subsubsection{Pre-Questionnaire}

\begin{enumerate}[label=\textbf{P\arabic*}, leftmargin=1.5cm, itemsep=0.8em]
\item Kommt auf die Anzahl Prüfungen an. Unter der woche etwa 3-4h und in den Ferien ca. 12-15h.

\item Shooter, Souls-likes, horror, action, cozy farming sim.

\item 1. visuelle coherenz und coole kunst, 2. animierende kämpfe mit vielen verschiedenen moves und kombos, 3. coole waffen, 4. character spezifische Eigenschaften.

\item Wort wörtlich also dass der Charakter Hackt und Slashed, also seine Gegner

\item 3-4

\item aggressiv und schnell

\item 1. So die grundrisse eines echten Spiels, also nur das nötigste, 2. Einen spielbaren Charakter mit flüssigen moves, 3. Anfänge eines world building, so man bekommt eine Idee vom konzept, 4. Bossfight und kampfspass.

\item Zu schwere gegner mit zu punishing moves wie z.b. marlenia.
\end{enumerate}

\subsubsection{Post-Questionnaire}

\paragraph{Gegnerverhalten}
\begin{enumerate}[label=\textbf{Q\arabic*}, leftmargin=1.5cm, itemsep=0.8em]
\item 5
\item 4
\item 5
\item 3
\end{enumerate}

\paragraph{Gegnerplatzierung und Kampfraum}
\begin{enumerate}[label=\textbf{Q\arabic*}, leftmargin=1.5cm, itemsep=0.8em, resume]
\item 5
\item 5
\item 3
\item 5
\end{enumerate}

\paragraph{Gesamteindruck}
\begin{enumerate}[label=\textbf{Q\arabic*}, leftmargin=1.5cm, itemsep=0.8em, resume]
\item 4
\item 5
\item 4
\item Die Verteilung der Gegner gefiel mir, also wie sie platziert/gespawnt wurden, deren leben war auch angemessen also kein unnötiges leben.
\item es war immer schwere wenn sie in einem haufen gewesen sind. Besonders wenn es die Soldiers waren. Springen unangenehm, weil ich erwartet habe das ich so einen krassen hit gegen den boden mach und gegner treffe. Kommt glaube von Elden Ring
\item Ich habe einen Boss erwartet, welcher nicht gekommen ist. Der Umfang der Spielmechaniken haben mich sehr überrascht. Der Spamgegner also der Soldier war angenehm zum spielen. Ich habe mir ein paar mehr sceen elemente vorgestellt, aslo mehr verzierung
\end{enumerate}

\subsubsection{Beobachtungsraster}

\paragraph{Gegnerverhalten}
\begin{enumerate}[label=\textbf{B\arabic*}, leftmargin=1.5cm, itemsep=0.8em]
\item
    \begin{itemize}[nosep]
        \item Die Testperson spielt insgesamt sehr aggressiv, ausser wenn es zu viele Gegner werden (siehe B2).
        \item Beim Heavy Soldier geht sie methodischer vor und versucht, keine Treffer zu nehmen.
        \item Die Heavy Soldiers tankt sie weiterhin einfach. Sie spielt für diesen Gegner zu aggressiv. Vielleicht ist der Wechsel zwischen aggressivem und taktischem Spielstil für sie eine zu starke Umstellung (Vermutung der Beobachtenden).
    \end{itemize}

\item
    \begin{itemize}[nosep]
        \item Sie sucht Abstand, wenn es zu viele Gegner werden, und fokussiert sich dann neu auf die Fernkämpfer.
        \item Bei zu vielen Gegnern rennt sie weg, um eine Atempause zu machen und zu heilen. Danach geht es aber weiter.
        \item Sie hat sich mittlerweile an Gruppen von vier bis fünf Gegnern gewöhnt.
        \item Sie kann mittlerweile drei Heavy Soldiers gleichzeitig bekämpfen.
    \end{itemize}

\item
    \begin{itemize}[nosep]
        \item Sie versucht meistens zuerst, den Archer auszuschalten. Da aber oft Soldiers im Weg stehen, hat sie hin und wieder Probleme.
        \item Bei niedriger Gesundheit ist ihr primäres Ziel, die Fernkämpfer auszuschalten.
        \item Wird sie von einem Pfeil getroffen, verschiebt sich ihr Fokus auf die Quelle des Schusses.
        \item Sobald sie von irgendeinem Archer getroffen wird, richtet sich ihr Fokus auf alle Fernkämpfer.
    \end{itemize}

\item
    \begin{itemize}[nosep]
        \item Die Testperson ignoriert eingehende Treffer, hält sie einfach aus und schlägt weiter.
        \item Sie tankt die Treffer weiterhin.
        \item Sie reagiert auf Angriffe erst, nachdem diese sie getroffen haben.
        \item Mittlerweile versucht sie, dem Heavy Archer auszuweichen, aber nur leicht.
    \end{itemize}

\item
    \begin{itemize}[nosep]
        \item Sie konnte seit Beginn nicht richtig springen, weil sie die Leertaste spammt. Das bereitet ihr weiterhin Schwierigkeiten.
        \item Sie versucht wiederholt, sich an einer Kante festzuhalten (Ledge Grab), zielt dabei aber auf eine Wand statt auf die greifbare Box und scheitert deshalb.
        \item Mittlerweile kann sie etwas besser springen.
    \end{itemize}
\end{enumerate}

\paragraph{Gegnerplatzierung und Kampfraum}
\begin{enumerate}[label=\textbf{B\arabic*}, leftmargin=1.5cm, itemsep=0.8em, resume]
\item
    \begin{itemize}[nosep]
        \item Bei neuen Encounters geht sie direkt auf den ersten Gegner zu, den sie finden kann.
        \item Zwischen zwei Räumen oder Wellen sucht sie eine Box zum Draufklettern, um den Raum zu überblicken und alle Gegner zu erfassen.
    \end{itemize}

\item
    \begin{itemize}[nosep]
        \item In offenen Encountern ist sie weniger vorsichtig und geht aggressiver vor. In engeren Encountern scheint sie eher strategisch zu spielen.
        \item In kleinen Räumen verhält sie sich etwas taktischer und weniger aggressiv, möglicherweise weil sie dort weniger Leben hatte. Sie teilt einen Treffer aus und zieht sich wieder zurück, bis alle Gegner tot sind (guerillaartig).
        \item Im kleinen Raum gibt es wenig Platz zum Heilen (siehe B12).
    \end{itemize}

\item
    \begin{itemize}[nosep]
        \item Bevor sie in den nächsten Raum geht, macht sie Pausen, um sich zu erholen und zu heilen.
        \item Zum Heilen geht sie an einen Ort ohne Gegner und nimmt Abstand.
        \item Sie benutzt die Boxen, um daraufzusteigen, kleine Pausen zu machen und zu heilen.
    \end{itemize}

\item
    \begin{itemize}[nosep]
        \item Sie geht direkt auf die ersten Heavy Soldiers zu und wird hart getroffen.
        \item Beim ersten Treffer durch einen Heavy Soldier ist sie überrascht. Sie wurde nur getroffen, weil sie ihn \enquote{trollen} wollte.
        \item Beim ersten Heavy Archer erkennt sie sofort, dass dessen Geschosse zielsuchend sind und sie verfolgen.
        \item Als der Chain Archer eingeführt wird, ist sie unsicher, ob das alle Gegner im Raum sind, und analysiert den gesamten Raum, bevor sie ihn angreift.
        \item Bei neuen Gegnern hält sie meistens kurz inne, bevor sie angreift.
    \end{itemize}
\end{enumerate}

\paragraph{Allgemein}
\begin{enumerate}[label=\textbf{B\arabic*}, leftmargin=1.5cm, itemsep=0.8em, resume]
\item
    \begin{itemize}[nosep]
        \item Nachdem der Encounter im Raum beendet war, bemerkt sie nicht gleich, dass sie weitergehen kann. Auch öffnende Türen bemerkt sie nicht immer sofort.
        \item Wenn sie keinen Gegner sieht, prüft sie die Tür. Ist diese geschlossen, sucht sie den Raum erneut ab.
    \end{itemize}

\item
    \begin{itemize}[nosep]
        \item Keine Einträge notiert.
    \end{itemize}

\item
    \begin{itemize}[nosep]
        \item Im kleinen Raum ist sie sichtlich genervt davon, dass sie beim Heilen unterbrochen wird.
        \item Besonders die Heavy Archers scheinen sie zu nerven (Eindruck der Beobachtenden).
    \end{itemize}

\item
    \begin{itemize}[nosep]
        \item Sie setzt manchmal unnötig Air Combat ein und trifft dann nicht alle Treffer.
        \item Sie verwendet selten Combo-Variationen, überwiegend nur die leichte Attacke.
    \end{itemize}

\item
    \begin{itemize}[nosep]
        \item Beim Wechsel von Raum zu Raum bemerkt sie nicht immer sofort, dass sich die Tür geöffnet hat (siehe B10). Davor erholt sie sich in der Regel (siehe B8).
    \end{itemize}
\end{enumerate}

\subsubsection{Semistrukturiertes Interview}

\begin{enumerate}[label=\textbf{I\arabic*}, leftmargin=1.5cm, itemsep=1em]
\item
\begin{enumerate}[label=\alph*), itemsep=0.6em]
    \item Du bist manchmal gegen die Wand gelaufen und hast versucht, dich daran zu greifen. Es hat öfters nicht funktioniert. Was hast du dir dabei überlegt? \\
        \textit{Antwort:} Ich hatte den Jump nicht ganz gecheckt, und ich wurde manchmal ungeduldig beim Springen und habe Ledge Grab versucht. Ich habe aber nicht ganz verstanden, dass ich gegen die Wand muss, und dass ich einfach Pech hatte, und habe es einfach noch einmal versucht. \\
        \textit{Bezug zur Beobachtung (B5):} Wiederholte Ledge-Grab-Versuche gegen die Wand statt gegen die greifbare Box. Sprünge gelangen anfangs schlecht und später etwas besser.
    \item Hast du die neuen Gegner immer nur anhand der Farbe erkannt, und konntest du an der Farbe den Typ erkennen? \\
        \textit{Antwort:} Ja, auch an den Attacken und an den Waffen, also z.\,B. an den verschiedenen Pfeilen. Und das Colorscheme, dass Archers so gelblich sind, war verständlich. \\
        \textit{Bezug zur Beobachtung (B9):} Beim ersten Heavy Archer erkannte sie sofort die zielsuchenden Geschosse. Bei neuen Gegnern hielt sie meist kurz inne.
\end{enumerate}

\item Du hattest zwei starke Reaktionen: einmal, als du Low Health hattest und viele Gegner da waren, und einmal, als du nicht hochgekommen bist. Was ist dir dabei durch den Kopf gegangen? \\
\textit{Antwort:} Ich dachte eigentlich, dass ich mal sterben will, um zu sehen, was passiert, und dann habe ich durchgezogen. Beim Nicht-Hochkommen dachte ich mir \enquote{Cooked, scheisse, dass ich es nicht geschafft habe}, und dann habe ich mich resettet, es nochmal versucht und geschafft. \\
\textit{Bezug zur Beobachtung (B2, B5):} Bei zu vielen Gegnern rannte sie weg, um zu heilen, und machte danach weiter. Beim Hochkommen scheiterte sie wiederholt am Ledge Grab.

\item An dieser Stelle bist du wiederholt gestorben bzw. hast du gezögert. Was hat dir das Weiterkommen erschwert? \\
\textit{Antwort:} Die Testperson ist nicht gestorben, die Frage entfällt.

\item Wie haben sich die grossen und kleinen Räume für dich unterschieden? \\
\textit{Antwort:} Mein Lieblingsraum war Raum~6, weil er eine schöne neue Struktur eingeführt hat, nämlich die Brücke. Damit kann man mehr machen als mit nur zwei Containern. Es gab eine coole Abwechslung zwischen grossen und kleinen Räumen, oder besser gesagt eine Mischung. Kleine Räume waren eher schwerer, weil es keine Zuflucht gab. \\
\textit{Bezug zur Beobachtung (B7, B12):} In kleinen Räumen spielte sie taktischer (Treffer austeilen und zurückziehen) und war genervt, wenn sie beim Heilen unterbrochen wurde.

\item Als hier mehrere Gegner gleichzeitig da waren: Wie hast du entschieden, wen du zuerst angreifst? \\
\textit{Antwort:} Ich bin fast immer zuerst auf Archers gegangen, wenn ich im Nahkampf war. Vor allem die, die drei Geschosse oder suchende Pfeile geschossen haben, waren krass, deswegen wollte ich die immer zuerst loswerden. Die mit den Hämmern konnte man gut countern, und sie waren langsam, deswegen habe ich sie zum Schluss aufbewahrt. Und die Soldiers haben mich vom Heilen abgehalten, weil sie schnell waren, und deswegen habe ich sie nach den Archers getötet. \\
\textit{Bezug zur Beobachtung (B3):} Sie griff meistens zuerst die Archers an, bei niedriger Gesundheit und nach Pfeiltreffern verstärkt. Oft standen ihr dabei Soldiers im Weg.

\item Empfandest du diesen Treffer bzw. diesen Tod als eigenen Fehler oder als unfair? Warum? \\
\textit{Antwort:} Wenn ich in der Luft getroffen wurde, war es deprimierend, bis ich wieder aufgestanden bin. War aber witzig. Homing Missiles waren schwerer zu dodgen, vor allem wenn man sie im letzten Moment bemerkt hat, und sie haben viel Schaden gemacht. Aber die Schüsse fühlten sich nicht fehl am Platz an und waren eher meine Fehler. Wenn viele Soldiers auf einmal gespawnt sind, war es ein wenig schwerer, lol, aber nicht unangenehm. \\
\textit{Bezug zur Beobachtung (B4, B13):} Sie reagierte auf Angriffe meist erst nach dem Treffer und tankte viele Treffer. Sie setzte manchmal unnötig Air Combat ein.

\item Gab es einen Punkt, an dem sich der Kampf für dich in einen anderen \enquote{Modus} verschoben hat (z.\,B. von entspannt zu angespannt)? Wo genau? \\
\textit{Antwort:} Ich fand, dass Heilpausen schwer zu finden waren, und ich musste einfach wegrennen, um das machen zu können. Die Waves kamen mehr und mehr, und man musste sich einen Moment erarbeiten oder auf Kisten stehen. Es fühlte sich an wie überlappende Wellen. Wenn die Räume fertig waren, konnte man entspannen, aber in einem Raum gab es keine allzu guten Möglichkeiten. Ausser im Trap Room, dort wusste ich, dass es einer werden würde. \\
\textit{Bezug zur Beobachtung (B8, B14):} Sie heilte an Orten ohne Gegner oder auf Boxen und erholte sich vor dem Raumwechsel. Beim Raumwechsel bemerkte sie nicht immer gleich, dass sich die Tür geöffnet hatte.

\item Gibt es etwas, das du während des Spielens gedacht oder gefühlt hast, das wir bisher nicht angesprochen haben? \\
\textit{Antwort:} Nichts, das mir jetzt einfallen würde. Also wir hatten kein Tutorial, und du hast es mir erklären müssen. Das wäre noch ein Verbesserungsvorschlag. \\
\textit{Bezug zur Beobachtung (B10):} Sie bemerkte nach Abschluss eines Raums nicht gleich, dass sie weitergehen konnte.
\end{enumerate}

\subsection{Testperson 3 (TP3)}

\subsubsection{Pre-Questionnaire}

\begin{enumerate}[label=\textbf{P\arabic*}, leftmargin=1.5cm, itemsep=0.8em]
\item 0 - 20h (kommt auf releases von Spielen an)

\item Action, Platformer

\item gutes Feedback; Dynamik; gut abgetrennte / erkennbare Segmente von Angriff und Verteidigung

\item Action hauptsächlich basierend auf Kampf; Spielzyklus mit ständig auftretenden neuen Erweiterungen; Steigerung der Fähigkeiten inner- und außerhalb des Spiels

\item 3

\item meist mit wenig Denken oder Konzipierung eines Spielstils, ich spiele wie ich mich fühle.

\item demo-artiges Schülerprojekt (bereits ein Teil gesehen)

\item leichteres Design gefällt tendenziell besser; aber keine besondere Präferenz oder Abneigung
\end{enumerate}

\subsubsection{Post-Questionnaire}

\paragraph{Gegnerverhalten}
\begin{enumerate}[label=\textbf{Q\arabic*}, leftmargin=1.5cm, itemsep=0.8em]
\item 5
\item 5
\item 5
\item 5
\end{enumerate}

\paragraph{Gegnerplatzierung und Kampfraum}
\begin{enumerate}[label=\textbf{Q\arabic*}, leftmargin=1.5cm, itemsep=0.8em, resume]
\item 4
\item 3
\item 5
\item 5
\end{enumerate}

\paragraph{Gesamteindruck}
\begin{enumerate}[label=\textbf{Q\arabic*}, leftmargin=1.5cm, itemsep=0.8em, resume]
\item 4
\item 1 für raum zu raum und 4 für locked in [= ins Spiel vertieft]
\item 5
\item Das taktische wen angreifen, healen, weglaufen, homing missiles kan man cool ausweichen. Der schadensunterschied zwischen links und rechts click dass man den archer mit rechtsclick instantly killen kann
\item die bewegungen sind schwerfällig und auch wenn ich das telegraphing sehe kann ich nicht immer reagieren. die range konnte ich nicht von anfang an einschätzen und es hatte eine learning curve.
\item Die bewegung hätte ich ein wenig flüssiger erwartet, ansonsten hat es ziemlich viel mit dem vorwissen und von dem was ich mitbekommen habe übereingestimmt.
\end{enumerate}

\subsubsection{Beobachtungsraster}

\paragraph{Gegnerverhalten}
\begin{enumerate}[label=\textbf{B\arabic*}, leftmargin=1.5cm, itemsep=0.8em]
\item
\begin{itemize}[nosep]
    \item Die Testperson ist ziemlich aggressiv, bewegt sich recht schnell und kommt zügig voran.
    \item Sie ignoriert ihr Low Health.
\end{itemize}

\item
\begin{itemize}[nosep]
    \item Zum ersten Mal rennt sie vor einer Dreiergruppe Soldiers weg, um zu heilen.
    \item Sie gerät bei vielen Gegnern nicht in Panik, sondern nur, wenn sie wenig Leben hat und viele Gegner da sind.
\end{itemize}

\item
\begin{itemize}[nosep]
    \item Momentan geht sie einfach von Archer zu Archer, lässt die Soldiers sie verfolgen und schlägt diese nebenbei mit.
    \item Den normalen Soldier empfindet sie als schwerer bzw. nerviger als den Heavy Soldier, weil der Heavy Soldier weniger oft trifft. Die Heavy Soldiers haben deshalb die kleinste Priorität.
    \item Sie greift oft an, was sie gerade vor sich sieht (siehe B6).
\end{itemize}

\item
\begin{itemize}[nosep]
    \item Sie versucht schon von Anfang an, Pfeilen auszuweichen, hält aber meistens einfach drauf und ignoriert den Schaden.
    \item Sie reagiert momentan selten auf Angriffe, besonders auf die der Nahkämpfer. Erst wenn nur noch Archers übrig sind, beginnt sie auszuweichen.
    \item Sie will zunehmend mehr ausweichen, kann es aber nicht gut, weil ihre eigenen Combos zu schwerfällig sind und ihr die Zeit zum Ausweichen fehlt.
    \item Sie ist generell gut in der Luftbewegung (Aerial Movement) und weicht dort oft aus.
\end{itemize}

\item
\begin{itemize}[nosep]
    \item Sie wird wiederholt in der Luft von Archers abgeschossen und fällt zu Boden.
    \item Sie springt weiterhin während des Kampfes, wahrscheinlich aus Gewohnheit (Muscle Memory) aus anderen Spielen. In diesem Spiel bringt das nichts.
    \item Das Ledge Grab beherrscht sie solide. Sie muss nicht pixelgenau springen, um auf die Box zu kommen, weil sie es kann und es weiss.
    \item Zur Luftkombo siehe B13.
\end{itemize}
\end{enumerate}

\paragraph{Gegnerplatzierung und Kampfraum}
\begin{enumerate}[label=\textbf{B\arabic*}, leftmargin=1.5cm, itemsep=0.8em, resume]
\item
\begin{itemize}[nosep]
    \item Sie greift oft den Gegner an, den sie gerade vor sich sieht.
    \item Einmal wird sie in der Luft von einem Archer getroffen, fällt zu Boden und wird dort von Soldiers und Heavy Soldiers umzingelt. Sie wird dafür deutlich bestraft.
    \item Sie findet einen versteckten Heavy Archer nicht und sucht ihn. Der Heavy Archer hat sie auch nicht aus dem Hinterhalt getroffen, deshalb hat sie ihn nicht bemerkt.
\end{itemize}

\item
\begin{itemize}[nosep]
    \item In grossen sowie in kleinen Räumen ist sie genau gleich aggressiv.
    \item Zu einem anderen Zeitpunkt ist sie in kleinen Räumen hektischer als in grossen und versucht, die Gegner schneller zu töten.
\end{itemize}

\item
\begin{itemize}[nosep]
    \item Sie stirbt einmal, weil sie nie geheilt hat.
    \item Zum ersten Mal heilt sie.
    \item Zwischen den Räumen macht sie kaum Pausen zum Heilen und geht einfach weiter.
    \item In Ruhephasen heilt sie inzwischen regelmässiger.
    \item Sie nimmt gern Abstand, um zu heilen. Das hat sie inzwischen gelernt (Eindruck der Beobachtenden).
    \item Sie hält eine ruhige Phase für das Ende des Raums (siehe B11).
    \item Zur Fluchtreaktion statt Box siehe B13.
\end{itemize}

\item
\begin{itemize}[nosep]
    \item Der Heavy Soldier interessiert sie nicht, sie greift ihn einfach an. Sie hat vorher gesagt, dass sie lieber ohne viel Nachdenken spielt (vgl. P6).
    \item Beim ersten Heavy Archer mit Homing Missile versucht sie auszuweichen, schafft es nicht und tankt den Treffer. Es scheint ihr egal zu sein.
    \item Den Chain Archer sieht sie zum ersten Mal. Es interessiert sie wenig, sie erkennt aber, dass es eine neue Art Gegner ist.
\end{itemize}
\end{enumerate}

\paragraph{Allgemein}
\begin{enumerate}[label=\textbf{B\arabic*}, leftmargin=1.5cm, itemsep=0.8em, resume]
\item
\begin{itemize}[nosep]
    \item Etwa in der Mitte des Spiels benutzt sie zum ersten Mal richtig den Double Jump. Es verwirrt sie, dass sie lange drücken kann, um höher zu springen, und sie nutzt es deshalb nicht. Sie spammt die Leertaste, möglicherweise weil sie etwas gestresst ist (Vermutung der Beobachtenden).
    \item Sie will in den nächsten Raum, kann aber nicht, und sucht deshalb den restlichen Gegner.
    \item Sie sucht den versteckten Heavy Archer (siehe B6).
\end{itemize}

\item
\begin{itemize}[nosep]
    \item Sie sagt, sie glaube, es sei die letzte Welle des Raums, obwohl sie wenig Leben hat. Es folgt aber noch eine weitere Welle. Weil es so ruhig war, dachte sie, es sei die letzte gewesen, hat sich aber getäuscht.
    \item Im zweitletzten Raum sagt sie \enquote{Boss} (laut Interview, siehe I2).
\end{itemize}

\item
\begin{itemize}[nosep]
    \item Sie wirkt sehr fokussiert.
    \item Sie hat sichtlich Spass und lacht die ganze Zeit.
    \item Besonders der Heavy Archer nervt sie.
    \item Einmal wird sie von Gegnern umzingelt und getroffen und ist sichtlich genervt (siehe I2).
    \item Nach dem Ende des Spiels spielt sie es direkt von vorn.
\end{itemize}

\item
\begin{itemize}[nosep]
    \item Sie versucht oft die Luftkombo, obwohl sie damit die Bodengegner seltener trifft.
    \item Statt auf eine Box zu steigen und Pause zu machen, versucht sie wegzurennen. Die Beobachtenden verstehen das nicht ganz, vielleicht hat sie es vergessen.
    \item Ihre momentane Taktik: Gegner anhäufen, hineinspringen, etwas Schaden machen und wieder herausgehen, immer wieder.
\end{itemize}

\item
\begin{itemize}[nosep]
    \item Sie macht zwischen den Räumen kaum Pausen und geht einfach weiter (siehe B8).
    \item Sie will in den nächsten Raum, kann aber nicht, und sucht deshalb den restlichen Gegner (siehe B10).
\end{itemize}
\end{enumerate}

\subsubsection{Semistrukturiertes Interview}

\begin{enumerate}[label=\textbf{I\arabic*}, leftmargin=1.5cm, itemsep=1em]
\item
\begin{enumerate}[label=\alph*), itemsep=0.6em]
\item Was hast du gedacht, bevor du zum ersten Mal gestorben bist? \\
    \textit{Antwort:} Ich hatte Heal vergessen und Leben ignoriert. \\
    \textit{Bezug zur Beobachtung (B1, B8):} Sie ignorierte ihr Low Health und starb, weil sie nie geheilt hatte.
\item Wie hast du die neuen Gegner wahrgenommen? \\
    \textit{Antwort:} Beim ersten Archer dachte ich, es ist einfach eine stärkere Version, und die anderen Archers habe ich auch nur am Angriff erkannt. Sonst habe ich die anderen schon bemerkt, wusste aber nicht genau, was sie machen. \\
    \textit{Bezug zur Beobachtung (B9):} Beim Chain Archer zeigte sie wenig Interesse, erkannte aber, dass es eine neue Art Gegner ist.
\end{enumerate}

\item
\begin{enumerate}[label=\alph*), itemsep=0.6em]
\item Warum hast du im zweitletzten Raum \enquote{Boss} gesagt? \\
\textit{Antwort:} Ja halt, grosser Raum und leer = Boss. \\
\textit{Bezug zur Beobachtung (B11):} Verbale Reaktion im zweitletzten Raum.
\item Du warst einmal von Gegnern umgeben, wurdest getroffen und warst sichtlich genervt. Warum? \\
\textit{Antwort:} Ich war stuck, es gab gefühlt einen Input Delay, und ich habe auf die Gegner nicht geachtet, und es hat mich genervt. \\
\textit{Bezug zur Beobachtung (B4, B12):} Ihre eigenen Combos waren zu schwerfällig, um rechtzeitig auszuweichen. Sichtbare Frustration beim Umzingeltwerden.
\end{enumerate}

\item An dieser Stelle bist du wiederholt gestorben bzw. hast du gezögert. Was hat dir das Weiterkommen erschwert? \\
\textit{Antwort:} Die Testperson ist nicht wiederholt gestorben, die Frage entfällt.

\item Wie unterscheiden sich die kleinen, mittleren und grossen Räume für dich? \\
\textit{Antwort:} Bei den kleinen Räumen wusste ich direkt, wohin. Bei den Räumen mit den Boxen hatte ich besonders viele Optionen, und der Raum mit der Brücke hatte am meisten Optionen. Gestresst war ich nicht, weil ich kontrollieren konnte, ob ich rein- oder rausgehe. \\
\textit{Bezug zur Beobachtung (B7):} Sie war in grossen und kleinen Räumen gleich aggressiv, zu einem anderen Zeitpunkt in kleinen Räumen hektischer.

\item Als hier mehrere Gegner gleichzeitig da waren: Wie hast du entschieden, wen du zuerst angreifst? \\
\textit{Antwort:} Definitiv erst mal die Archers, vor allem die normalen, danach die anderen Archer-Typen. Die waren halt immer omnipräsent, und ich habe viel Schaden kassiert. Die Nahkampfgegner kommen sowieso zu einem, und es gibt keinen Unterschied zwischen gross und klein. Die kleineren hatten nicht wirklich Präsenz. Die grossen hatten Momente, in denen man dachte, ich muss ausweichen, aber die haben auch oft nicht getroffen. \\
\textit{Bezug zur Beobachtung (B3):} Sie ging von Archer zu Archer und liess die Soldiers folgen. Heavy Soldiers hatten die kleinste Priorität.

\item Empfandest du diesen Treffer bzw. diesen Tod als eigenen Fehler oder als unfair? Warum? \\
\textit{Antwort:} Ja, nach den Attack-Combos sieht man zum Beispiel, wie Arrows kommen, und man kann trotzdem nicht ausweichen, weil der Input nicht freigegeben wurde. Aber man hatte genug Healing-Möglichkeiten. \\
\textit{Bezug zur Beobachtung (B4, B5):} Ihre eigenen Combos waren zu schwerfällig, um rechtzeitig auszuweichen. Sie wurde wiederholt in der Luft abgeschossen.

\item Gab es einen Punkt, an dem sich der Kampf für dich in einen anderen \enquote{Modus} verschoben hat (z.\,B. von entspannt zu angespannt)? Wo genau? \\
\textit{Antwort:} Nachdem man alle Ranged Enemies getötet hat, war es immer entspannt. Davor ist es recht hektisch, und ich muss aufpassen. Aber es ergab eine gute Balance von Spannung und Entspannung. Am Anfang Spannung, am Ende Entspannung. Bugs zu finden hat mir besonders Freude bereitet. \\
\textit{Bezug zur Beobachtung (B8, B11, B14):} In Ruhephasen heilte sie inzwischen regelmässiger. Sie hielt eine ruhige Phase fälschlich für das Raumende. Zwischen den Räumen machte sie kaum Pausen.

\item Wie hast du das Erscheinen der Gegner und ihre Platzierung wahrgenommen? \\
\textit{Antwort:} Die Gegner sind zwar ausserhalb meiner Sichtweite gespawnt, man hat sie aber noch früh genug bemerkt. Ich glaube, die Gegnerplatzierung ist stellenweise festgesetzt, und danach spawnen halt Archers eher erhöht und in Gruppen. Also die erste Welle ist definitiv immer gesetzt, zumindest teilweise. \\
\textit{Bezug zur Beobachtung (B6):} Sie griff meist den Gegner vor sich an. Einen versteckten Heavy Archer fand sie nicht.
\end{enumerate}

\subsection{Testperson 4 (TP4)}

\subsubsection{Pre-Questionnaire}

\begin{enumerate}[label=\textbf{P\arabic*}, leftmargin=1.5cm, itemsep=0.5em]
\item Früher mehrere Stunden, heutzutage ist es ziemlich genau 0h. :(
\item Am liebsten sind Taktik-Spiele oder Zivilisationsaufbauspiele, mit Kriegselementen. Jedoch ist auch ab und zu ein Spiel, welches sich nur auf Geballer konzentriert, nett.
\item Gute Visualisierung (Sprich, Gegner muss sichtbar sein, das Umfeld soll schön sein und zur atmosphäre beitragen…), ein smoothes Handling und befriedigende Combos.
\item Unter einem HAck and Slash Spiel, stelle ich mir etwas vor, dass nicht viel Hirnleistung erfordert und einfach da ist um cool auszusehen, bzw. dich gut fühlen zu lassen mit spektakulären combos etc. pp.
\item 3
\item Beides hat seine Reize, wenn ich müde bin, dann zweiteres, sonst möchte ich Hirnfutter.
\item Ein mediocre af game, was aber spass macht und an manchen stellen unnötig krass gut sein wird, an anderen wiederum den anschein erweckt, 3 Tage lang gemacht worden zu sein.
\item Unsichtbare Feinde ohne jegliche Andeutung machen keinen Sinn. Und einfach 3 Minuten einem Gegner ausweichen zu müssen, ohne eine Gegenspiel Möglichkeit zu haben ist auch wack. Zu schwierige Gegener sind probably auch nervig.
\end{enumerate}

\subsubsection{Post-Questionnaire}

\paragraph{Gegnerverhalten}
\begin{enumerate}[label=\textbf{Q\arabic*}, leftmargin=1.5cm, itemsep=0.5em]
\item 5
\item 4
\item 3
\item 1
\end{enumerate}

\paragraph{Gegnerplatzierung und Kampfraum}
\begin{enumerate}[label=\textbf{Q\arabic*}, leftmargin=1.5cm, itemsep=0.5em, resume]
\item 2
\item 2
\item 3
\item 5
\end{enumerate}

\paragraph{Gesamteindruck}
\begin{enumerate}[label=\textbf{Q\arabic*}, leftmargin=1.5cm, itemsep=0.5em, resume]
\item 4
\item 4
\item 5
\item Horden an kleinen Viechern zu zerstören. Das movement war sehr cool gemacht, hatte aber keine wirklichen Auswirkungen
\item Die Combos haben sehr gestört, da nie ganz klar war wann ich wieder angreifen kann oder verwundbar bin, zudem konnten sie nicht abgebrochen werden.
\item Ähnlich zum erwartetem, wobei bei der Variation der Gegner positiv überrascht, blocken und weitere Kontrollierbare Elemente hätte ich aber schon erwartet.
\end{enumerate}

\subsubsection{Beobachtungsraster}

\paragraph{Gegnerverhalten}
\begin{enumerate}[label=\textbf{B\arabic*}, leftmargin=1.5cm, itemsep=0.8em]
\item
\begin{itemize}[nosep]
\item Die Testperson ist die ganze Zeit ungefähr auf halbem Leben.
\item Nachdem sie gestorben ist, geht sie schneller durch die Räume und spielt deutlich aggressiver.
\item Sie spielt weniger taktisch und greift eher einfach an.
\end{itemize}
\item
\begin{itemize}[nosep]
\item Sie häuft Soldiers an, bevor sie diese angreift, und beginnt anzugreifen, wenn sie bei einem Archer steht, um alles auf einmal treffen zu können.
\item Nach einem Treffer wird sie beim Aufstehen erneut getroffen und kann nichts machen.
\end{itemize}
\item
\begin{itemize}[nosep]
\item Sie schlägt zuerst die Gegner, die auf sie zukommen.
\item Nachdem ein Heavy Archer sie getroffen hat, konzentriert sie sich auf die Heavy Archers.
\item Sie schlägt weiterhin auf den ersten Gegner, der vor ihr erscheint.
\item Mittlerweile geht sie zuerst auf die Archers.
\item Sie geht eher auf die Soldiers als auf die Heavy Soldiers, weil die Wahrscheinlichkeit kleiner ist, dass die Heavy Soldiers sie treffen.
\end{itemize}
\item
\begin{itemize}[nosep]
\item Sie weicht den Projektilen aus und hält Abstand zu den Archers.
\item Sie reagiert nicht auf Schaden, sondern läuft einfach hinein und schlägt.
\item Sie dasht nicht viel.
\item Sie versucht auszuweichen, wenn es geht, läuft aber meistens einfach durch.
\end{itemize}
\item
\begin{itemize}[nosep]
\item Sie spammt die Dash-Taste, statt sie gedrückt zu halten.
\end{itemize}
\end{enumerate}

\paragraph{Gegnerplatzierung und Kampfraum}
\begin{enumerate}[label=\textbf{B\arabic*}, leftmargin=1.5cm, itemsep=0.8em, resume]
\item
\begin{itemize}[nosep]
\item Sobald sie eine Welle sieht, schaut sie sich diese zuerst an, bevor sie angreift.
\item Sie stirbt in Raum~12, weil sie einen Archer nicht gesehen hat und von ihm getroffen wurde.
\item Sie springt und positioniert sich in der Luft, bevor sie in die Gegnerwelle eintaucht.
\end{itemize}
\item
\begin{itemize}[nosep]
\item Keine Einträge notiert.
\end{itemize}
\item
\begin{itemize}[nosep]
\item Sie heilt sich, bevor sie in den nächsten Raum geht.
\end{itemize}
\item
\begin{itemize}[nosep]
\item Sie schlägt ohne nachzudenken auf den neu eingeführten Archer ein.
\item Beim Heavy Archer erkennt sie sofort, dass es ein neuer Gegner ist.
\item Sie bemerkt, dass der Chain Archer neu ist, und greift ihn nicht an, um zu sehen, was er macht.
\end{itemize}
\end{enumerate}

\paragraph{Allgemein}
\begin{enumerate}[label=\textbf{B\arabic*}, leftmargin=1.5cm, itemsep=0.8em, resume]
\item
\begin{itemize}[nosep]
\item Sie bemerkt nicht, dass die Tür offen ist.
\item Sie bewegt die Kamera hektisch hin und her.
\item Sie bemerkt die offene Tür erneut nicht.
\item Bei einem Angriff bemerkt sie nicht, dass sie einen Gegner getroffen hat.
\item Sie kann nicht genau einschätzen, wie lange das Heilen dauert.
\end{itemize}
\item
\begin{itemize}[nosep]
\item Sie meint, die Gegner reagierten irgendwie nicht auf Treffer. Als sie genauer hinsieht, fällt es ihr aber auf.
\item Sie meint, eine Jump-Drop-Attacke wäre cool.
\item Sie sagt, sie geniesse ihre Überlegenheit.
\item Sie sagt, es mache eigentlich schon noch ganz Spass.
\item Sie meint, es sei besonders befriedigend, wenn man alle auf einmal besiegen kann.
\item Sie meint, das Angreifen sei etwas clunky, das sei für unseren Combat aber nicht so schlimm.
\item Sie meint, das Einzige, was im Spiel bestrafe, seien die etwas längeren Kampfanimationen und Treffer in der Luft.
\item Sie benutzt lieber die Light Attack, weil die Heavy Attack ihrer Meinung nach für ihre Dauer zu wenig Schaden macht.
\item Sie findet besonders den Heavy Archer schlimm, weil sie von ihm getroffen wird. Die Chain Archers stören sie nicht.
\item Sie würde die Combos lieber abbrechen können, versteht aber, wieso sie so designt sind.
\item Sie findet es nicht cool, von hinten angegriffen zu werden.
\item Sie findet es spassig, die Gegner vorherzusagen, weil sie die Räume schon gesehen hat.
\item Sie findet die Ästhetik schön.
\item Sie wünscht sich einen Block.
\item Sie findet das Movement richtig cool und meint, es habe eine Lernkurve.
\item Sie findet die Animationen des Heavy Soldiers und des Chain Archers cool.
\end{itemize}
\item
\begin{itemize}[nosep]
\item Sie ist erstaunt, als sie in den dritten Raum kommt.
\end{itemize}
\item
\begin{itemize}[nosep]
\item Sie benutzt die Flying Combo ohne Grund.
\item In Raum~1 ignoriert sie die Gegner und probiert das Movement aus.
\item Sie benutzt vor allem die Light-Attack-Combo.
\item Sie versucht, mit Objekten zu interagieren.
\item Sie benutzt gerade zum ersten Mal die Heavy Attack.
\item Sie springt gerne in Gegnermengen hinein, weil sie meinte, es gebe eine Sprungattacke (vgl. B11).
\item Sie schaut beim Schlagen kaum noch hin, weil sie ohnehin trifft.
\end{itemize}
\item
\begin{itemize}[nosep]
\item Beim Raumwechsel bemerkt sie zweimal nicht, dass die Tür offen ist (siehe B10). Vor dem Raumwechsel heilt sie sich (siehe B8).
\end{itemize}
\end{enumerate}

\subsubsection{Semistrukturiertes Interview}

\begin{enumerate}[label=\textbf{I\arabic*}, leftmargin=1.5cm, itemsep=1em]
\item
\begin{enumerate}[label=\alph*), itemsep=0.6em]
\item Was ist dir durch den Kopf gegangen, als du zum zweiten Mal gestorben bist? \\
\textit{Antwort:} Ja, war halt dumm gelaufen mit den Combos, welche einfach ned aufgehört haben.
\item Was ist dir bei den Air Combos durch den Kopf gegangen? \\
\textit{Antwort:} Useless af, es gab halt dann einfach nur nervige Situationen, wo ich in der Luft gehangen bin.
\end{enumerate}

\item Wie hast du den Heavy Soldier wahrgenommen? \\
\textit{Antwort:} Ja, ist ein bisschen ein Sponge, was Schaden angeht, war aber mehr lustig, dass er die Hälfte der Zeit verfehlt hat und sonst nur in Gruppen zur Gefahr wurde, was bissl konträr ist.

\item Was ist dir beim Drop-in durch den Kopf gegangen? \\
\textit{Antwort:} Ja, ned ganz gecheckt, dass ich eigentlich nicht gepunisht werde für grosse Höhen. Ansonsten war halt grosse Arena, aber nichts los ausser einem Lörach.

\item Wie haben sich die beiden Zwischenräume, der grosse und der kleine, für dich unterschieden? \\
\textit{Antwort:} Für mich hat es keinen Unterschied gemacht, vllt war ich aber auch einfach nur im Bann von Hauen und Smashen.

\item Als hier mehrere Gegner gleichzeitig da waren: Wie hast du entschieden, wen du zuerst angreifst? \\
\textit{Antwort:} Ja, wenn ich wollte, Archer, weil die halt mies nerven und alle unterschiedlich angreifen, ohne wirkliche Prediction oder Soundeffekte, die mich warnen könnten. Sonst auf YOLO.

\item Empfandest du diesen Treffer bzw. diesen Tod als eigenen Fehler oder als unfair? Warum? \\
\textit{Antwort:} Naja, irgendwo Skill Issue, war aber auch blöd, da die Combos manchmal mies gebackfired haben und ich nicht blocken konnte oder so shyyt.

\item Gab es einen Punkt, an dem sich der Kampf für dich in einen anderen \enquote{Modus} verschoben hat (z.\,B. von entspannt zu angespannt)? Wo genau? \\
\textit{Antwort:} Ich war eigentlich die ganze Zeit ins Spiel vertieft und habe mich nur darauf konzentriert, möglichst schnell die Maus zu drücken, bin ich ehrlich, [Name].

\item Gibt es etwas, das du während des Spielens gedacht oder gefühlt hast, das wir bisher nicht angesprochen haben? \\
\textit{Antwort:} Ja, meine geheimen Bestechungsgelder wurden nicht empfangen. NA FR ON GOD JETZT: Ich habe das Gefühl, dass man mehr Interaktionen hätte bauen sollen zum Spielen, wobei Button Mashing schon funny ist, aber währenddessen auch noch Block oder Dodge Roll noch funnier wäre. YFM?
\end{enumerate}

\subsection{Testperson 5 (TP5)}

\subsubsection{Pre-Questionnaire}

\begin{enumerate}[label=\textbf{P\arabic*}, leftmargin=1.5cm, itemsep=0.5em]
\item je nach dem 4-30h
\item Singerlplayer open World, Factory games, multiplayer sandbox games, soulslike, rouglikes
\item challenging, gutes design, eigenes( gute controlles sei es movement, aim, parry), gut passende musikalische untermalung und soundeffekts
\item sehr viele gegner töten
\item In action games 4
\item methodisch, alles mögliche ausprobieren
\item Ich erwarte ein fertige fighting game mit einstimmer grafik gutem gamedesigne und das es spass macht, Gute setting für anpassungen
\item unfaire treffer/gegner, allgemein dinge die man nicht vorhersehen kann, dann aber dafür bestraft wird.
\end{enumerate}

\subsubsection{Post-Questionnaire}

\paragraph{Gegnerverhalten}
\begin{enumerate}[label=\textbf{Q\arabic*}, leftmargin=1.5cm, itemsep=0.5em]
\item 3
\item 4
\item 4
\item 3
\end{enumerate}

\paragraph{Gegnerplatzierung und Kampfraum}
\begin{enumerate}[label=\textbf{Q\arabic*}, leftmargin=1.5cm, itemsep=0.5em, resume]
\item 2
\item 5
\item 5
\item 5
\end{enumerate}

\paragraph{Gesamteindruck}
\begin{enumerate}[label=\textbf{Q\arabic*}, leftmargin=1.5cm, itemsep=0.5em, resume]
\item 4
\item 4
\item 5
\item Wenn du jemanden Spawnkillst, also wenn du schon weisst wo er spawnt oder du vermutest wo dieser spawn.
\item Wenn ich stunlocked worden bin von ganz vielen gegnern.
\item ich dachte ich wäere noch viel stärker also wie hyrule warriors und ich habe erwartet dass ich 1000de umbringe
\end{enumerate}

\subsubsection{Beobachtungsraster}

\paragraph{Gegnerverhalten}
\begin{enumerate}[label=\textbf{B\arabic*}, leftmargin=1.5cm, itemsep=0.8em]
\item
\begin{itemize}[nosep]
\item Die Testperson geht, anders als die anderen Testpersonen, eher taktisch als aggressiv vor.
\item Sie geht inzwischen ziemlich schnell durch die Räume, aber immer noch eher taktisch.
\end{itemize}
\item
\begin{itemize}[nosep]
\item Sie will die Gegner gruppieren und dann besiegen, kassiert dabei aber viel Schaden.
\item Sie rennt im Kreis, um sich Taktiken zu überlegen. Währenddessen trifft sie nichts.
\item Sie sammelt die Gegner bei einem Archer und versucht, alle zu besiegen, schafft es aber nicht immer ganz sauber.
\item Sie versucht jetzt, Gegner abzusondern und einzeln anzugreifen, weil ihr Gruppenkämpfe zu viel Schaden machen. Das Absondern ist aber nicht so einfach.
\end{itemize}
\item
\begin{itemize}[nosep]
\item Manchmal ignoriert sie die Soldiers und geht direkt auf die Archers, manchmal geht sie zuerst auf die Nahkämpfer.
\item Sie findet die Archers nervig, weil diese sie treffen.
\item Sie kämpft vor allem mit Soldiers und greift nur hin und wieder einen Archer an, wenn dieser sie abgeschossen hat.
\item Sie greift zuerst die Gegner auf den Boxen an, um sich eine Zuflucht zu sichern, falls es knapp wird.
\end{itemize}
\item
\begin{itemize}[nosep]
\item Sie findet die Ankündigungsanimation der Archers etwas verwirrend, weil es aussieht, als würden sie laufen, obwohl sie stehen bleiben.
\item Wenn sie Schaden nimmt, zieht sie sich zurück, wo andere Testpersonen einfach weiter draufgehalten hätten. Das bestraft sie mehr: Haben alle Gegner einmal angegriffen, kann man sie leicht besiegen. Sie zieht sich aber zurück und kann deshalb keinen Schaden machen (Eindruck der Beobachtenden).
\item Sie ist daher eher empfindlich gegenüber Schaden (Eindruck der Beobachtenden).
\end{itemize}
\item
\begin{itemize}[nosep]
\item Keine Einträge notiert.
\end{itemize}
\end{enumerate}

\paragraph{Gegnerplatzierung und Kampfraum}
\begin{enumerate}[label=\textbf{B\arabic*}, leftmargin=1.5cm, itemsep=0.8em, resume]
\item
\begin{itemize}[nosep]
\item Sie greift an, was vor ihr ist.
\item Sie wird beinahe in einen Türrahmen gedrängt und kommt wegen der Gegnermengen fast nicht mehr heraus.
\end{itemize}
\item
\begin{itemize}[nosep]
\item In Raum~10 ist sie vorsichtig, wenn sie um die Ecken geht.
\end{itemize}
\item
\begin{itemize}[nosep]
\item Sie stellt sich auf Props, um vor den Gegnern Zuflucht zu finden.
\item Sie heilt sich, wenn am Ende einer Welle nur noch wenige Gegner übrig sind.
\item Sie sucht zum Heilen erhöhte Positionen (High Ground).
\item Sie stirbt, weil sie zu wenig Leben hat und mangels Mana nicht mehr heilen kann.
\end{itemize}
\item
\begin{itemize}[nosep]
\item Bei neuen Gegnern greift sie an, ohne zu schauen, was diese machen.
\item Sie greift auch den Heavy Soldier ohne zu zögern an.
\end{itemize}
\end{enumerate}

\paragraph{Allgemein}
\begin{enumerate}[label=\textbf{B\arabic*}, leftmargin=1.5cm, itemsep=0.8em, resume]
\item
\begin{itemize}[nosep]
\item Sie kann nicht ganz einschätzen, wie gross ihre Reichweite ist.
\item Sie hat etwas Mühe mit dem Wechsel zwischen Movement und Combat.
\end{itemize}
\item
\begin{itemize}[nosep]
\item Sie möchte emoten können.
\item Sie findet, es spiele sich wie \game{Hyrule Warriors}, das sie gespielt hat und das zum selben Genre gehört.
\item Sie möchte Heavy- und Light-Attack-Tasten zu Combos kombinieren können. Im Prototyp gibt es aber nur Combos mit Pausen.
\item Sie findet, der Heavy Soldier springe einfach irgendwohin.
\item Sie meint, sie könne nicht einschätzen, wie viel Schaden eine Gruppe von Gegnern macht, wenn sie hineinläuft.
\item Sie meint, sie erkenne die verschiedenen und neuen Gegner an Farbe und Aussehen (vgl. B9).
\item Sie meint, sie sei nie unter Stress gewesen.
\item Sie meint, der Hammer-Sound passe irgendwie nicht.
\end{itemize}
\item
\begin{itemize}[nosep]
\item Sie wirkt sehr selbstsicher.
\end{itemize}
\item
\begin{itemize}[nosep]
\item Sie versucht, Glitches zu finden.
\item Sie springt zuerst einmal die Klippe hinunter, um zu sehen, was passiert.
\item Sie versucht, durch die Map zu clippen.
\item Sie prüft, ob Gegner Pfeile abfangen können.
\item Sie versucht, die Gegner-KI zu durchschauen.
\item Sie sucht nach interaktiven Objekten. Es gibt aber keine.
\item Sie wechselt zwischen Heavy und Light Attack ab.
\item Sie versucht, die Gegner von ihrem High Ground aus auszutricksen, was aber nicht funktioniert.
\item Sie bemerkt, dass die Angriffe unterschiedlich viel Schaden machen und dass zwei Heavy Attacks einen Archer direkt besiegen.
\item Sie ist die ganze Zeit eher damit beschäftigt, das Spiel zu beobachten und Muster sowie Bugs zu finden.
\end{itemize}
\item
\begin{itemize}[nosep]
\item Sie hat keine Probleme zu erkennen, dass die Türen offen sind.
\item In Raum~11 steht sie direkt über dem Drop-in und gelangt sofort in Raum~12. Sie erschrickt kurz, überblickt dann aber direkt den Raum und beruhigt sich.
\end{itemize}
\end{enumerate}

\subsubsection{Semistrukturiertes Interview}

\begin{enumerate}[label=\textbf{I\arabic*}, leftmargin=1.5cm, itemsep=1em]
\item Du bist nach einem Raum manchmal einfach stehen geblieben und hast nichts gemacht. \\
\textit{Antwort:} Hab die Map angesehen, den nächsten Raum, die Musik und nach Bugs gesucht.

\item Du hast hier eine Reaktion gezeigt bzw.\ etwas gesagt. Kannst du beschreiben, warum? \\
\textit{Antwort:} --

\item Woran haben deine Tode meistens gelegen? Du bist ja dreimal gestorben. \\
\textit{Antwort:} Ich weiss nicht, passiert halt. Ich hab manchmal overcommittet und habe Hits kassiert und habe dann nicht gepunisht.

\item Wie fühlen sich kleine Räume anders an als grosse? \\
\textit{Antwort:} Kleine sind übersichtlicher, es gibt aber die Gefahr, in eine Ecke gedrängt zu werden. Die grossen haben halt mehr Gegner, was sie für mich schwerer macht.

\item Als hier mehrere Gegner gleichzeitig da waren: Wie hast du entschieden, wen du zuerst angreifst? \\
\textit{Antwort:} Bogenschützen zuerst und danach den Rest. Und hab gleichzeitig versucht zu groupen.

\item Empfandest du diesen Treffer bzw. diesen Tod als eigenen Fehler oder als unfair? Warum? \\
\textit{Antwort:} Tode waren sehr punishing, weil du alles neu machen musstest, aber ich empfand sie nicht als unfair, weil ich sie kommen gesehen habe. Manchmal konnte ich aber nicht ganz einschätzen, wie viel Damage ich kassiere, wenn ich jetzt eine Gruppe angreife.

\item Gab es einen Punkt, an dem sich der Kampf für dich in einen anderen \enquote{Modus} verschoben hat (z.\,B. von entspannt zu angespannt)? Wo genau? \\
\textit{Antwort:} Hat hauptsächlich damit zu tun gehabt, wie viel Leben ich habe. Wenig Leben = Stress.

\item Gibt es etwas, das du während des Spielens gedacht oder gefühlt hast, das wir bisher nicht angesprochen haben? \\
\textit{Antwort:} Das Design finde ich geil, und den Artstyle habe ich gefühlt. Ich hätte gerne noch ein Extra-Feature gehabt, dass sich ein Balken auflädt für eine Special Attack. Fliegende Gegner wären cool. Von der Luft abgeschossen zu werden, war sehr punishing. Ein Checkpoint wäre auch noch geil gewesen. Combos waren nicht so intuitiv, weil ich mehr Mixes gewohnt bin, nicht Rhythmen. Das Beste am Spiel war wirklich der Combat und das Movement.
\end{enumerate}

\newpage
\section{Anhang: Glossar} \label{sec:Glossar}

In diesem Glossar erklären wir Fachbegriffe, die wir in dieser Arbeit verwenden.\footnote{Mit Claude formuliert, siehe KI-Deklaration (Abschnitt~\ref{sec:KIDeklaration}).}

\begin{description}[style=unboxed, leftmargin=0cm, itemsep=0.3em]
\item[Air Combo] Combo, die in der Luft gegen einen hochgeschlagenen Gegner ausgeführt wird.
\item[Asset] Vorgefertigter Bestandteil eines Spiels, etwa ein 3D-Modell, eine Textur, eine Animation, Musik oder ein Soundeffekt (Abschnitt~\ref{sec:Assets}).
\item[Beat 'em up] Genre, in dem sich die Spielfigur im Nahkampf durch viele Gegner kämpft. Vorläufer des Hack-and-Slash.
\item[Bossgegner] Besonders starker, meist einzigartiger Gegner, häufig am Ende eines Levels.
\item[Bullet-Hell] Spielprinzip, bei dem sehr viele gegnerische Projektile gleichzeitig in Mustern über den Bildschirm fliegen, denen der Spieler ausweichen muss.
\item[Combo] Abfolge von Angriffen, die durch eine bestimmte Tastenfolge nacheinander ausgelöst wird. In \game{Project Polearm} gibt es eine Light, eine Heavy und eine Air Combo (Tabelle~\ref{tab:Combos}).
\item[Crunch] Phase mit stark erhöhtem Arbeitsaufwand, oft mit Überstunden, kurz vor einer Deadline. In der Spieleentwicklung verbreitet.
\item[Dash] Kurze, schnelle Ausweich- oder Sprintbewegung.
\item[Encounter] Einzelne Kampfbegegnung. In unserem Prototyp findet jeder Encounter in einem Raum statt und besteht aus mehreren Wellen. Encounter~3 und~12 finden im selben Raum statt (Abschnitt~\ref{sec:Kampfabschnitt}).
\item[Feedback] Rückmeldung des Spiels auf eine Eingabe oder ein Ereignis, etwa durch Animation, Ton, Kamerabewegung oder Anzeigen.
\item[Frame] Einzelbild. Spiele berechnen meist 30 bis 60 oder mehr Bilder pro Sekunde. Zeitangaben bei Animationen werden oft in Frames gemacht.
\item[Game Engine] Software-Grundgerüst für die Spieleentwicklung (Darstellung, Physik, Eingabe, Audio). Wir verwenden Godot.
\item[Greybox] Rohfassung eines Levels aus einfachen, einfarbigen Grundformen, mit der Raumaufteilung und Spielbarkeit getestet werden, bevor Grafik hinzukommt.
\item[Hack-and-Slash] Genre der Action-Spiele mit Fokus auf Echtzeit-Nahkampf mit Waffen (Abschnitt~\ref{sec:Genredefinition}).
\item[Hitstun] Kurze Zeitspanne nach einem Treffer, in der die getroffene Figur nicht handeln kann.
\item[I-Frames] Kurz für \textit{Invincibility Frames}: kurzes Zeitfenster, meist während eines Ausweichmanövers, in dem die Spielfigur nicht getroffen werden kann.
\item[Input-Buffer] Zwischenspeicher für Eingaben, die etwas zu früh erfolgen. Sie werden ausgeführt, sobald die laufende Aktion es zulässt.
\item[Intimate Space] Raumtyp nach Totten: kleiner, überschaubarer Raum (Abschnitt~\ref{sec:Raumtypologien}).
\item[Knockback] Rückstoss, der eine getroffene Figur ein Stück wegschleudert.
\item[Ledge Grab] Festhalten an einer Kante, um sich daran hochzuziehen.
\item[Level] Abgeschlossener Spielabschnitt bzw.\ die Spielumgebung, in der er stattfindet.
\item[Mana] Ressource, die in \game{Project Polearm} zum Heilen verbraucht wird.
\item[Metaprogression] Dauerhafter Fortschritt über einzelne Level hinweg, etwa durch freigeschaltete Fähigkeiten oder Gegenstände.
\item[Narrow Space] Raumtyp nach Totten: enger, schmaler Raum, der die Bewegung einschränkt (Abschnitt~\ref{sec:Raumtypologien}).
\item[Pacing] Steuerung von Tempo und Intensität über den Spielverlauf, etwa durch den Wechsel von Anspannung und Erholung (Abschnitt~\ref{sec:Pacing}).
\item[Playtest] Test, bei dem Testpersonen ein Spiel spielen und dabei beobachtet und befragt werden (Abschnitt~\ref{sec:VorgehenPlaytest}).
\item[Pre-/Post-Questionnaire] Fragebogen vor bzw.\ nach dem Playtest.
\item[Prospect Space] Raumtyp nach Totten: offener, grosser Raum mit Überblick (Abschnitt~\ref{sec:Raumtypologien}).
\item[Rig] Digitales Skelett eines 3D-Modells, mit dem es animiert wird.
\item[Scaffolding] Schrittweises Heranführen an neue Spielelemente, zuerst einzeln, dann in Kombination (Abschnitt~\ref{sec:Scaffolding}).
\item[Spannungskurve] Verlauf der Intensität über einen Spielabschnitt.
\item[Spawn, Spawnpunkt] Das Erscheinen einer Figur im Level bzw.\ der Ort, an dem sie erscheint.
\item[Stylemeter] Anzeige in Stylish-Action-Spielen wie \game{Devil May Cry}, die bewertet, wie abwechslungsreich und elegant der Spieler kämpft.
\item[Stylish-Action] Spiele des Hack-and-Slash-Genres, die elegantes und abwechslungsreiches Kämpfen ins Zentrum stellen (Abschnitt~\ref{sec:Genredefinition}).
\item[Telegrafieren] Ankündigen eines gegnerischen Angriffs durch sichtbare oder hörbare Signale, damit der Spieler reagieren kann (Abschnitt~\ref{sec:V4}).
\item[Think-aloud] Methode, bei der Testpersonen ihre Gedanken während des Spielens laut äussern (Abschnitt~\ref{sec:MethodePlaytest}).
\item[Third-Person-Perspektive] Kameraperspektive hinter bzw.\ über der Spielfigur, sodass die Figur selbst sichtbar ist.
\item[Welle] Gruppe von Gegnern, die gemeinsam erscheint. Die nächste Welle folgt erst, wenn die vorherige besiegt ist (Abschnitt~\ref{sec:Kampfabschnitt}).
\end{description}

\newpage
\section{Anhang: Liste der externen Assets} \label{sec:Assets}

Die folgende Liste führt die externen Ressourcen auf, die im Prototyp verwendet oder als Grundlage angepasst wurden. Der Einsatz von KI ist in der KI-Deklaration (Abschnitt~\ref{sec:KIDeklaration}) aufgeführt. Quellenangaben zu übernommenem Code finden sich in der \textit{Project Polearm Technical Documentation}.

\paragraph{3D-Modelle und Animationen}
\begin{itemize}[leftmargin=1.5cm]
\item \textbf{Y Bot} (Grundmodell für Spielfigur und Gegner) und Animationen -- Mixamo (Adobe), \url{https://www.mixamo.com}. Lizenz: Nutzungsbedingungen von Mixamo. Angepasst: Farbvarianten für die Gegner, Animationen in Blender überarbeitet.
\item \textbf{Spear Animations Pack} -- 9CG, Fab, \url{https://www.fab.com/listings/02b4b496-4591-4cee-8970-e28c08026eae}, Fab Standard License. Alle Animationen in Blender abgeändert.
\item \textbf{Fantasy Treasure Chest} (Truhe) -- -X-ScornGames, Sketchfab, \url{https://sketchfab.com/3d-models/x-fantasy-treasure-chest-313a236caaec44f5b55aed327d7b56d8}, CC BY 4.0. Laut Angabe des Autors mit KI generiert.
\item \textbf{Umgebungstexturen und Props} (u.\,a.\ Bodenfliesen, Fels, Beton, Metallplatten, Rohre, Lüftungsschächte, Feuerlöscher, Feuermelder, Generator, Stromkasten, Überwachungskamera, Kanaldeckel) -- Poly Haven, \url{https://polyhaven.com}, CC0.
\item \textbf{Prototype Textures 1.0} (Greybox-Texturen) -- Kenney, \url{https://www.kenney.nl}, CC0.
\end{itemize}

\paragraph{Shader}
\begin{itemize}[leftmargin=1.5cm]
\item \textbf{Stylized Sky with Procedural Sun and Moon} -- krzmig, godotshaders.com, \url{https://godotshaders.com/shader/stylized-sky-with-procedural-sun-and-moon/}, MIT-Lizenz.
\item \textbf{Godot Foam Water} (Wasser-Shader) -- antzGames, GitHub, \url{https://github.com/antzgames/godot-foam-water}, MIT-Lizenz.
\end{itemize}

\paragraph{Benutzeroberfläche}
\begin{itemize}[leftmargin=1.5cm]
\item \textbf{UI-Elemente} (Panels, Rahmen, Trennlinien) -- Kenney, \url{https://www.kenney.nl}, CC0.
\item \textbf{Tastensymbole} (Tastatur und Maus, Xbox, PlayStation) -- Kenney, \url{https://www.kenney.nl}, CC0.
\item \textbf{Schrift Noto Serif SC} -- Google Fonts, SIL Open Font License 1.1.
\end{itemize}

\paragraph{Musik}
\begin{itemize}[leftmargin=1.5cm]
\item \textbf{Fight Them Until We Can't} -- Zander Noriega, OpenGameArt, \url{https://opengameart.org/content/fight-them-until-we-cant}, CC-BY 3.0. Verwendung: Spiel.
\item \textbf{Big Boss} -- bart, OpenGameArt, \url{https://opengameart.org/content/big-boss}, CC-BY 3.0. Verwendung: Spiel.
\item \textbf{Dark Ambience Loop} -- Iwan Gabovitch (qubodup), OpenGameArt, \url{https://opengameart.org/content/dark-ambience-loop}, CC-BY 3.0 (wahlweise auch GPL 2.0 oder 3.0, von uns unter CC-BY 3.0 verwendet). Verwendung: Pausenmusik.
\item \textbf{Theme of Com-Mecha} (\enquote{Welcome to Com-Mecha}) -- Matthew Pablo, OpenGameArt, \url{https://opengameart.org/content/theme-of-com-mecha}, CC-BY 3.0. Verwendung: Hauptmenü.
\item \textbf{Boss Battle \#2 [Symphonic Metal]} -- nene, OpenGameArt, \url{https://opengameart.org/content/boss-battle-2-symphonic-metal}, CC0. Verwendung: Hauptmenü.
\item \textbf{Bleeding Out} -- Brandon Morris, OpenGameArt, \url{https://opengameart.org/content/bleeding-out}, CC0 bzw.\ OGA-BY 3.0. Verwendung: Todesbildschirm.
\end{itemize}

\paragraph{Soundeffekte}
\begin{itemize}[leftmargin=1.5cm]
\item \textbf{Sword Combat Sound Effects Pack (Free Version)} -- Hove Audio, itch.io, \url{https://hoveaudio.itch.io/sword-combat-sound-effects-pack-free-version}. Lizenz: lizenzfrei (Royalty Free), eine Nennung ist laut Autor nicht nötig.
\item \textbf{400 Sounds Pack} -- Chequered Ink, itch.io, \url{https://ci.itch.io/400-sounds-pack}. Lizenz: freie Nutzung, auch kommerziell, mit oder ohne Nennung, solange die Dateien nicht unverändert weiterverkauft werden.
\item \textbf{Metal Hit} -- Marcelo Fernandez, OpenGameArt, \url{https://opengameart.org/content/metal-hit}, CC-BY 3.0.
\end{itemize}

\paragraph{Godot-Erweiterungen}
\begin{itemize}[leftmargin=1.5cm]
\item \textbf{FuncGodot} (Import der TrenchBroom-Karten) -- func\_godot Community, MIT-Lizenz.
\item \textbf{Phantom Camera} -- Marcus Skov, MIT-Lizenz.
\item \textbf{Debug Draw 3D} -- Dmitriy Salnikov, MIT-Lizenz.
\item \textbf{Animation Player Refactor} -- poohcom1, MIT-Lizenz.
\end{itemize}

\end{document}
